% Round 16. Author: model_obstruction. Independent review required.
% A new criterion relative to the frozen archive, not a priority claim.
\section{Finite symmetries of the canonical numerical residual}
\label{r16sym:section}

The residual of Round 15 is canonical enough to carry the symmetries
of a supplied model.  Its total isotropy can then force vanishing
without a sign condition on the primitive intersection form.  This
section does not prove that an arbitrary variety has such a model
or such symmetries.

\subsection{The exact geometric input}

Fix a finite extension $K/\mathbf Q_p$, its absolute value, its
valuation ring $R$, and its finite residue field $k$.  Let
$\mathcal X/R$ be regular, integral, flat, and projective of
relative dimension $n\geq4$, with smooth geometrically integral generic fiber
$X$.  Assume
\[
                  \mathcal X_k=\sum_{i\in I}m_iD_i
\tag{ER16.1}\label{r16sym:special-fiber}
\]
has strict normal crossings reduced support and smooth native
closed strata.  The $m_i$ are its scheme-theoretic multiplicities.
All numerical groups below have rational coefficients, and all
degrees are native degrees over $k$.

Let $H$ be one relatively ample line bundle, used on all strata,
and write
\[
\begin{gathered}
 V^a=\bigoplus_iN^a(D_i),\qquad
 A_1u=\left(j_i^!\sum_jj_{j*}u_j\right)_i,\\
 P_H^2=\bigoplus_i\ker\left(
     H^{n-3}:N^2(D_i)\longrightarrow N^{n-1}(D_i)\right),\\
 Q_{H,m}(x,y)=\sum_i m_i\deg_{D_i}(x_iy_iH^{n-4}).
\end{gathered}
\tag{ER16.2}\label{r16sym:numerical-data}
\]
For one prime $\ell\ne p$, put
\[
 C_X^2=\ker\left(\operatorname{cl}_\ell:
     \operatorname{CH}^2(X)_{\mathbf Q}
       \longrightarrow H^4(X_{\bar K},\mathbf Q_\ell(2))\right).
\]

We use exactly the following frozen results, with their original
model hypotheses and native normalizations.
Theorem \texttt{r15res:canonical-residual}, in
\path{round15/actual_local/canonical_primitive_numerical_residual.tex},
constructs a rational-linear map
\[
                    R_{\mathcal X,H}:C_X^2\longrightarrow P_H^2.
\tag{ER16.3}\label{r16sym:residual-map}
\]
If $\Theta$ extends $z$ in the model Chow group and
$T=(j_i^!\Theta)_i$, choose a rational solution of
$H^{n-3}A_1u=H^{n-3}T$ and set
$R_{\mathcal X,H}(z)=T-A_1u$.  The answer is independent
of both choices.  Its zero condition is equivalent to a
componentwise numerically zero model extension and implies
$\mathfrak o_X^2(z)=0$ and strong G on the prescribed exact
support complement.  Proposition \texttt{r15iso:total-isotropy},
in \path{round15/actual_local/weighted_primitive_isotropy.tex},
proves
\[
             Q_{H,m}(R_{\mathcal X,H}(z),R_{\mathcal X,H}(w))=0
                      \quad(z,w\in C_X^2).
\tag{ER16.4}\label{r16sym:frozen-isotropy}
\]
In particular, this is total isotropy of the whole image, not
only the vanishing of a norm for one chosen cycle.

\begin{definition}[Finite model symmetry]
\label{r16sym:datum}
A finite model symmetry is a finite group $G$ acting on
$\mathcal X$ by $R$-automorphisms and preserving the numerical
class of $H$.  We use the left action
$g\cdot z=(g^{-1})^*z$ on generic Chow groups and the analogous
permuted pullback on the component numerical groups.

A rational $G$-module is called multiplicity-free if each
simple $\mathbf Q[G]$-module occurs with multiplicity at most
one.  This is a statement over $\mathbf Q$, not over a
splitting field and not about the dimension of the invariant
subspace alone.
\end{definition}

If a finite group already acts on the model, a preserved ample
class can always be chosen by replacing an ample line bundle
$L$ with $\bigotimes_{g\in G}g^*L$.  The multiplicity-free
condition must be checked for the resulting primitive space;
we do not assert that it is independent of the polarization.

\begin{lemma}[Equivariance and isometry]
\label{r16sym:equivariance}
For a finite model symmetry, $P_H^2$ is a rational
$G$-module, $Q_{H,m}$ is $G$-invariant, and
\[
                  R_{\mathcal X,H}(g\cdot z)
                         =g\cdot R_{\mathcal X,H}(z)
                         \qquad(g\in G, z\in C_X^2).
\tag{ER16.5}\label{r16sym:equivariance-formula}
\]
Thus $W=\operatorname{im}R_{\mathcal X,H}$ is a
$G$-stable totally isotropic rational subspace of $P_H^2$.
\end{lemma}

\begin{proof}
An $R$-automorphism preserves the Cartier divisor of a
uniformizer.  It therefore permutes the $D_i$ and preserves
their scheme multiplicities $m_i$.  Isomorphic native
components and strata have the same native degree functional.
The assumed equality of numerical ample classes makes
multiplication by $H$ equivariant.  Hence the kernel defining
$P_H^2$ is stable and the sum defining $Q_{H,m}$ is invariant.

Restriction and regular Gysin commute with the Cartesian
squares defined by these isomorphisms.  Therefore $A_1$ and
$H^{n-3}A_1$ commute with the permuted pullback action.  Given
an extension $\Theta$ and a solution $u$ for $z$, the extension
$(g^{-1})^*\Theta$ and the vector $g\cdot u$ solve the same
normalization equation for $g\cdot z$.  Substituting them in
the defining formula for the residual proves
\eqref{r16sym:equivariance-formula}.  Geometric cycle-class
vanishing is preserved by these pullbacks, so the source
$C_X^2$ is itself stable.  Rational linearity and
\eqref{r16sym:frozen-isotropy} prove the last assertion.
No numerical-to-homological implication in codimension two
is used.
\end{proof}

\subsection{The representation-theoretic argument}

\begin{lemma}[A multiplicity-free orthogonal module]
\label{r16sym:linear-algebra}
Let $G$ be finite and $P$ a finite-dimensional rational
$G$-module.  Let $Q$ be a nondegenerate $G$-invariant symmetric
bilinear form on $P$.  If $P$ is multiplicity-free, its only
$G$-stable totally isotropic subspace is zero.

More generally, write $P$ as its rational isotypic direct sum
$\bigoplus_\rho P_\rho$.  Distinct $P_\rho$ are
$Q$-orthogonal and every restricted form $Q|_{P_\rho}$ is
nondegenerate.  The same assertions hold for any sum of
isotypic summands.
\end{lemma}

\begin{proof}
Maschke's averaging argument makes every rational finite-group
module semisimple.  Every rational simple module is self-dual:
choose a rational basis, average its positive definite
coordinate form over $G$, and use the resulting rational
$G$-invariant nondegenerate pairing to identify the module
with its dual.  This auxiliary positive definite form is a
linear-algebra construction; it makes no assertion about
the geometric intersection form $Q$.

The bilinear restriction $Q:P_\rho\times P_\sigma\to
\mathbf Q$ corresponds to a $G$-map
$P_\rho\to P_\sigma^*$.  For distinct rational simple
types it is zero by semisimplicity and Schur's lemma.  Hence
the isotypic decomposition is orthogonal.  Nondegeneracy of
$Q$ on $P$ then implies nondegeneracy on each summand and on
every sum of them.

Suppose now that $P$ is multiplicity-free and $W\subset P$
is stable.  Maschke's theorem supplies a stable complement
$U$.  Since each simple type occurs at most once in $P$,
$W$ and $U$ have no common simple type.  The self-duality
argument gives $Q(W,U)=0$.  If $W$ is totally isotropic,
also $Q(W,W)=0$.  Thus $W$ lies in the radical of $Q$ on
$P=W\oplus U$, which is zero.
\end{proof}

\begin{theorem}[Vanishing from finite model symmetries]
\label{r16sym:vanishing}
Under the model assumptions \eqref{r16sym:special-fiber},
suppose $G$ is a finite model symmetry, $Q_{H,m}$ is
nondegenerate on $P_H^2$, and $P_H^2$ is multiplicity-free
as a rational $G$-module.  Then
\[
                       R_{\mathcal X,H}=0.
\tag{ER16.6}\label{r16sym:zero}
\]
Every class in $C_X^2$ has a componentwise numerically zero
model extension.  Every specified rational cycle $Z$ in
such a class has an ordinary actual form
\[
             \gamma_Z\in B^{1,1}
                    (X^{\mathrm{an}}\setminus|Z|^{\mathrm{an}})
\]
satisfying
\[
 \int_{X^{\mathrm{an}}\setminus|Z|^{\mathrm{an}}}
                         \gamma_Z\wedge D\eta
                  =-\int_{Z^{\mathrm{an}}}\eta
                 \quad\text{for every }
                 \eta\in B_c^{n-2,n-2}(X^{\mathrm{an}}).
\tag{ER16.7}\label{r16sym:strong-G}
\]
\end{theorem}

\begin{proof}
The image of the residual is stable and totally isotropic
by Lemma~\ref{r16sym:equivariance}.  Lemma~\ref{r16sym:linear-algebra}
makes it zero.  The numerical and actual conclusions are
then exactly the zero-residual conclusions of the frozen
Theorem \texttt{r15res:canonical-residual}.  In particular,
they use the pure operational model witness, the original
field model-Chern detector, and the established sufficient
direction $\mathfrak o_X^2([Z])=0\Rightarrow G(Z)$.
No converse to that direction is required.
\end{proof}

\begin{corollary}[The fourfold criterion]
\label{r16sym:fourfold}
If $n=4$, nondegeneracy in
Theorem~\ref{r16sym:vanishing} is automatic.  Thus a finite
model symmetry acting multiplicity-free on $P_H^2$ implies
strong G for all geometrically $\ell$-homologically trivial
codimension-two cycles on $X$, with no upper bound on
$\dim_{\mathbf Q}P_H^2$.
\end{corollary}

\begin{proof}
On a smooth projective fourfold $D_i$, the middle numerical
pairing $N^2(D_i)\times N^2(D_i)\to\mathbf Q$ is
nondegenerate by the definition of numerical equivalence.
Its restriction to $HN^1(D_i)$ is the divisor pairing
$\deg(uvH^2)$, which is nondegenerate by the divisor Hodge
index theorem.  The numerical orthogonal complement of
$HN^1(D_i)$ is precisely
$\ker(H:N^2(D_i)\to N^3(D_i))$.  Consequently the primitive
restriction is nondegenerate.  The direct sum and the
nonzero weights $m_i$ preserve nondegeneracy.  This is also
the argument recorded in the proof of frozen Corollary
\texttt{r15iso:fourfold-rank-one}; no codimension-two Hodge
standard conjecture enters.
\end{proof}

\begin{corollary}[Selected rational isotypic cycle sectors]
\label{r16sym:sectors}
Let $n=4$ and let $e\in\mathbf Q[G]$ be a sum of rational
central primitive idempotents.  Suppose $eP_H^2$ is
multiplicity-free.  Then
\[
                 R_{\mathcal X,H}(ez)=0
                              \qquad(z\in C_X^2),
\tag{ER16.8}\label{r16sym:isotypic-vanishing}
\]
and every specified cycle representing a class in $eC_X^2$
satisfies the full conclusion \eqref{r16sym:strong-G}.
It is enough that the simple types selected by $e$ each
occur with multiplicity at most one; other isotypic
summands may have arbitrary multiplicity.
\end{corollary}

\begin{proof}
Equivariance gives $R(ez)=eR(z)$.  The space
$W_e=R(eC_X^2)$ is $G$-stable because $e$ is central and
is totally isotropic by \eqref{r16sym:frozen-isotropy}.
The fourfold argument makes $Q_{H,m}$ nondegenerate on
$P_H^2$, and the isotypic assertion of
Lemma~\ref{r16sym:linear-algebra} makes its restriction to
$eP_H^2$ nondegenerate.  That lemma now forces $W_e=0$.
The frozen exact-support bridge applies to every specified
representative, including a representative whose support
is unrelated to the individual orbit representatives
used to describe the class $ez$.
\end{proof}

\paragraph{Two directly checkable forms of the hypothesis.}
In relative dimension four it suffices that $P_H^2$ be an
irreducible rational representation of the finite model
symmetry group.  More generally, suppose commuting model
involutions act on $P_H^2$, preserve $H$, and each simultaneous
rational eigenspace has dimension at most one.  Their finite
elementary abelian group acts multiplicity-free, so the
same theorem applies.  For one involution, it is enough
that both primitive eigenspaces have dimension at most one.

\begin{proposition}[A cohomological certificate for the symmetry criterion]
\label{r16sym:cohomological-certificate}
In the fourfold setting, let $e\in\mathbf Q[G]$ be a sum of
rational central primitive idempotents.  If, for some
$\ell\ne p$, the finite-dimensional representation
\[
       e\left(\bigoplus_i
          H^4((D_i)_{\bar k},\mathbf Q_\ell(2))\right)
\tag{ER16.9}\label{r16sym:cohom-representation}
\]
is multiplicity-free as a $\mathbf Q_\ell[G]$-module,
then $eP_H^2$ is multiplicity-free over $\mathbf Q$.
Consequently every class in $eC_X^2$ satisfies strong G.
No Tate conjecture and no identification of numerical
and homological equivalence in codimension two is required.
\end{proposition}

\begin{proof}
For a native component $D_i$, let $A_{i,\ell}$ be the image
of the geometric cycle map from
$\operatorname{CH}^2(D_i)_{\mathbf Q}\otimes\mathbf Q_\ell$
to its displayed cohomology.  There is a natural surjection
\[
                  A_{i,\ell}\longrightarrow
                         N^2(D_i)_{\mathbf Q}\otimes\mathbf Q_\ell.
\tag{ER16.10}\label{r16sym:cycle-image-quotient}
\]
Indeed, if a finite $\mathbf Q_\ell$-linear combination of
cycle classes has zero cohomology class, its intersection
degree against every native complementary cycle is zero
by compatibility of the cycle map with intersection and
trace.  The perfect numerical pairing after scalar
extension makes its numerical class zero.  Thus the
numerical quotient factors through the cohomological
image, and it is surjective by the definition of the
numerical cycle group.  This direction is homological
vanishing implying numerical vanishing, not its converse.

The direct sum of \eqref{r16sym:cycle-image-quotient} is
$G$-equivariant because the action comes from model
isomorphisms, including the permutation of components.
The primitive numerical group is a $G$-stable subspace
of the target.  Hence $eP_H^2\otimes\mathbf Q_\ell$ is
a subquotient of \eqref{r16sym:cohom-representation}.
Maschke semisimplicity shows that a subquotient of a
multiplicity-free representation is multiplicity-free.

If a rational simple representation occurred twice in
$eP_H^2$, its nonzero scalar extension to $\mathbf Q_\ell$
would have at least one simple constituent, occurring
at least twice in $eP_H^2\otimes\mathbf Q_\ell$.
This contradicts the preceding conclusion.  Thus
$eP_H^2$ is multiplicity-free over $\mathbf Q$.
Apply Corollary~\ref{r16sym:sectors}.
\end{proof}

This certificate deliberately uses the full geometric
degree-four cohomology in the selected isotypic sector.
Replacing it by a smaller cohomological primitive group
would require a separate proof that it surjects onto the
numerical primitive sector; such a comparison is not
silently assumed.

\paragraph{Why this is a different sufficient criterion.}
The elementary example
$P=\mathbf Q^2$, $Q(x,x)=x_1^2-x_2^2$, with an involution
$(x_1,x_2)\mapsto(x_1,-x_2)$, is multiplicity-free and has
no nonzero invariant totally isotropic subspace, although
$Q$ has nonzero rational isotropic vectors.  Thus the new
linear-algebra criterion is not an implication of rational
anisotropy.  This example is only a linear-algebra comparison;
it is not asserted to be the primitive numerical lattice
of a fourfold, nor is it a counterexample to a Hodge
standard conjecture.  No new geometric family beyond the
stated, verifiable symmetry condition has been identified
in this section.

\subsection{Original-field dominant propagation}

\begin{corollary}[Invariant numerical ranks on a dominant model]
\label{r16sym:invariant-cover}
Let $X/K$ be smooth projective geometrically integral of
dimension at least two.  Suppose that after a finite
extension $F/K$ there is a projective dominant morphism
$f:V\to X_F$, where $V/F$ is smooth projective geometrically
integral of dimension $n\geq4$, and a supplied model
$\mathcal V/F^\circ$ with the same regular SNC hypotheses
and a finite model symmetry $G$ preserving $H$.
Assume $f\circ g=f$ on the generic fiber for every $g\in G$.
Put $V^a=\bigoplus_iN^a(D_i)$ for this model.

Either of the following conditions is sufficient for
$\mathfrak o_X^2$ to vanish on the complete geometric
homological kernel, and hence for strong G on every
prescribed cycle in that kernel:
\begin{enumerate}
\item $\dim (V^2)^G=\dim (V^1)^G$, in any relative
      dimension $n\geq4$;
\item $n=4$ and
      $\dim (V^2)^G-\dim (V^1)^G\leq1$.
\end{enumerate}
In particular a Galois alteration with a regular
$G$-equivariant model may be used when it satisfies the
stated invariant-rank condition.  Existence of alterations
with that condition is not asserted.
\end{corollary}

\begin{proof}
The frozen numerical divisor splitting
$V^2=H V^1\oplus P_H^2$ is $G$-equivariant.  Taking
invariants and using injectivity of multiplication by
$H$ gives
\[
             \dim(P_H^2)^G
                    =\dim(V^2)^G-\dim(V^1)^G.
\tag{ER16.11}\label{r16sym:invariant-rank}
\]
For a geometric homologically trivial class $z$ on $X$,
the pullback $f^*z$ is fixed by $G$.  Equivariance puts
its residual in $(P_H^2)^G$.

Under the first condition this space is zero.  Under
the second condition the fourfold pairing is
nondegenerate on $P_H^2$.  Its restriction to the
trivial isotypic summand $(P_H^2)^G$ is nondegenerate
by Lemma~\ref{r16sym:linear-algebra}.  A nondegenerate
rational form on a space of dimension at most one has
no nonzero isotropic vector, so the residual is again
zero by \eqref{r16sym:frozen-isotropy}.

The vanishing supplies a numerical model witness on
$V$.  Regard the model over $K^\circ$, retaining
ramification multiplicities and residue-field degrees,
and apply the frozen original-field numerical detector.
Then frozen dominant degree-two Bott--Chern injectivity
and the exact-support bridge give the assertion on $X$.
The same native-field argument is written in detail in
Corollary~\ref{r16sym:dominant} below.  It does not require
the alteration to be finite on any prescribed complement.
\end{proof}

\begin{corollary}[A dominant symmetric model witness]
\label{r16sym:dominant}
Let $X/K$ be smooth geometrically integral projective with $\dim X\geq2$.
Suppose there is a finite extension $F/K$, a smooth geometrically integral
projective $F$-variety $V$, and a projective dominant
$F$-morphism $f:V\to X_F$.  Suppose $V$ has a supplied
regular flat projective model over $F^\circ$ satisfying the
geometric and symmetry hypotheses of
Theorem~\ref{r16sym:vanishing}, or the fourfold criterion
of Corollary~\ref{r16sym:fourfold}.

Then every geometrically $\ell$-homologically trivial
codimension-two class on $X$ has zero actual obstruction.
Every specified representative satisfies strong G over
the original $K$, on its exact support complement, against
all original actual tests.
\end{corollary}

\begin{proof}
The pullback of a geometric homologically trivial class
to $V$ is again geometric homologically trivial.  The new
criterion gives a numerical model witness there.  Regard
that model as a model over $K^\circ$.  Its special-fiber
multiplicities acquire the ramification factor and its
supported degrees acquire the residue-field factor;
numerical vanishing is preserved.  The frozen native-field
model-witness argument therefore gives
$\mathfrak o_{V,K}^2(f^*z)=0$ directly in the original
$K$-actual theory.  There is no identification of different
test categories by a change of notation.

Apply frozen Theorem \texttt{r15dom:dominant} and Corollary
\texttt{r15dom:obstruction-witness}, in
\path{round15/root/general_degree_two_descent.tex}, to the
composite dominant $K$-morphism $V\to X_F\to X$.
Its finite maps retain their complete native function-field
degrees, including constant-field factors.  They descend
the global curvature primitive; the exact-support Green
pair is then formed for the original prescribed cycle.
No Green form is pushed along an arbitrary alteration.
\end{proof}

\paragraph{Audit boundaries.}
The new step is the equivariance--isotropy argument and its
cycle sectors.  The residual, weighted isotropy, model
detector, and dominant degree-two descent are frozen inputs.
For $n>4$, nondegeneracy of $Q_{H,m}$ remains an explicit
additional assumption; this text does not deduce it from
divisor Hodge index.  The representation condition must
hold for actual model automorphisms, not for an arbitrary
abstract action invented on the numerical vector space.
A nonzero residual never implies failure of strong G.
Neither general strong G nor a canonical unconditional
BB height on the complete geometric homological kernel
is concluded.

\paragraph{Literature check and status.}
The representation-theoretic input is elementary Maschke
semisimplicity and Schur orthogonality, for which a primary
author reference is J.~S.~Milne, \emph{Group Theory},
Section 7, Theorem 7.4 and the following semisimplicity
discussion, \url{https://www.jmilne.org/math/CourseNotes/GTe6.pdf}.
The argument is proved above in the exact rational scope used.
A focused literature search was performed during this round;
no claim of priority or comprehensive novelty is made.
